\subsection{HNL Construction and List Reuse}
\label{sec:hnl_measured_cost}

We first compare HNL with direct recursive MAC traversal from the source-tree
roots. Both methods use the same adaptive meshes, targets, MAC, list storage,
and OpenMP schedule; only the source search differs. The parameters are
$(p,P,\theta)=(6,4,0.5)$ and 16 CPU threads. Table~\ref{tab:hnl_ablation}
reports medians of three interleaved runs. HNL time includes construction of
the level-indexed vertex incidence, candidate generation, cover
materialization, and MAC classification. Validation is excluded from both
timings, although every run is checked for complete, nonoverlapping source
coverage and MAC compliance.

\begin{table}[!htbp]
\centering
\caption{HNL and direct root-recursive MAC construction. Visit counts are
averages per target; $M_{\rm HNL}$ is the number of cover-materialization
visits.}
\label{tab:hnl_ablation}
\small
\setlength{\tabcolsep}{4.2pt}
\begin{tabular}{rrrrrrr}
\toprule
$N$ & \multicolumn{2}{c}{$T_{\rm list}$ (s)} &
\multicolumn{2}{c}{MAC visits} & reduction & $M_{\rm HNL}$\\
\cmidrule(lr){2-3}\cmidrule(lr){4-5}
& HNL & root & HNL & root & (\%) &\\
\midrule
$10\,328$  & 0.127 & 0.060 & $1\,856$ & $1\,916$ & 3.15 & 507\\
$34\,464$  & 0.701 & 0.249 & $2\,801$ & $2\,901$ & 3.46 & 861\\
$100\,656$ & 2.075 & 0.723 & $2\,768$ & $2\,868$ & 3.47 & 856\\
$191\,544$ & 5.491 & 1.724 & $3\,549$ & $3\,683$ & 3.64 & $1\,235$\\
$822\,720$ & 25.131 & 7.735 & $3\,723$ & $3\,863$ & 3.62 & $1\,317$\\
\bottomrule
\end{tabular}
\end{table}

HNL reduces MAC visits by $3.1$--$3.6\%$, but candidate generation and
materialization prevent a reduction in total list-construction time. The near
lists are identical; the far-list counts differ by at most $0.62\%$ because
clusters may be accepted at different levels. On the $N=100\,656$ mesh, the
HNL and root partitions give absolute errors $2.1886\times10^{-4}$ and
$2.2434\times10^{-4}$, respectively. We therefore use HNL as the
topology-based construction in this implementation, without claiming a speed
advantage over direct root traversal or an asymptotic bound for its
construction cost.

The lists depend on the mesh, targets, and $\theta$, but not on the density.
We test reuse on the $N=100\,656$ mesh with 64 CPU threads using
\begin{equation}
V_k=V_0+\beta_k V_1,
\qquad \rho_k=-\frac{1}{4\pi}\Delta V_k,
\qquad \beta_k=0.5(0.8)^k,
\label{eq:hnl_reuse_sequence}
\end{equation}
where $V_0$ is given by Eq.~\eqref{eq:manufactured_solution} and
$V_1=\exp[-6\pi(x-0.70)^2-5\pi(y+0.35)^2-9\pi(z-0.20)^2]$.
One path rebuilds the lists at every update; the other retains the initial
lists and updates only density samples, moments, and the potential.

\begin{figure}[!htbp]
\centering
\resizebox{\textwidth}{!}{%
\begin{tikzpicture}
\begin{groupplot}[
  group style={group size=2 by 1, horizontal sep=0.55cm},
  journalaxis,
  width=0.48\textwidth,
  height=0.36\textwidth,
  xmin=0, xmax=21,
  xtick={1,2,5,10,20}
]
\nextgroupplot[
  xlabel={number of density updates},
  ylabel={cumulative time (s)},
  ymin=0, ymax=58,
  legend style={at={(0.03,0.97)},anchor=north west,draw=none}
]
\addplot+ coordinates {
  (1,2.798) (2,5.570) (5,13.684) (10,27.245) (20,53.934)
};
\addlegendentry{Rebuild}
\addplot+ coordinates {
  (1,2.950) (2,4.857) (5,10.612) (10,20.193) (20,38.988)
};
\addlegendentry{Reuse}

\nextgroupplot[
  xlabel={number of density updates},
  ylabel={speedup from reuse},
  ymin=0.95, ymax=1.42,
  ytick pos=right,yticklabel pos=right
]
\addplot+ coordinates {
  (1,0.991) (2,1.154) (5,1.290) (10,1.349) (20,1.381)
};
\addplot[dashed,black,mark=none] coordinates {(1,1) (20,1)};
\end{groupplot}
\end{tikzpicture}%
}
\caption{Cumulative time with list reconstruction and reuse (left), and the
resulting speedup (right). Values are medians of three runs.}
\label{fig:hnl_list_reuse}
\end{figure}

Reuse saves 25.9\% after ten updates and 27.6\% after twenty updates, giving
a twenty-update speedup of 1.381. This manufactured sequence demonstrates
amortization on a fixed mesh; it is not a complete Kohn--Sham SCF test.
